\QHeading{第7题　脱式计算。}
\par\noindent


\CalcBlock[简便解法（补充）]{1}{\begin{aligned}[t]&\phantom{={}}\color{black}768\div 12\times 24\\&\color{answerred}=768\times (24\div 12)\\&\color{answerred}=768\times 2\\&\color{answerred}=1536\end{aligned}}
\par
\MethodBlock{按运算顺序计算（基础练习）}{本单元还要练习运算顺序；与上面的简便解法对照。}{\begin{aligned}[t]&\phantom{={}}768\div12\times24\\&=64\times24\\&=1536\end{aligned}}


\CalcBlock{2}{\begin{aligned}[t]&\phantom{={}}\color{black}760-912\div24\times5\\&\color{answerred}=760-38\times5\\&\color{answerred}=760-190\\&\color{answerred}=570\end{aligned}}
\CalcBlock{3}{\begin{aligned}[t]&\phantom{={}}\color{black}817-93\div3+260\\&\color{answerred}=817-31+260\\&\color{answerred}=786+260\\&\color{answerred}=1046\end{aligned}}
\CalcBlock[方法一]{4}{\begin{aligned}[t]&\phantom{={}}\color{black}(85+79)\times(85-79)\\&\color{answerred}=164\times6\\&\color{answerred}=984\end{aligned}}
\MethodBlock{方法二：乘法分配律}{把括号中的两部分分别相乘，再相加或相减。}{\begin{aligned}[t]&\phantom{={}}(85+79)\times (85-79)\\
&=(85+79)\times 85-(85+79)\times 79\\
&=164\times 85-164\times 79\\
&=13940-12956\\
&=984\end{aligned}}

\CalcBlock{5}{\begin{aligned}[t]&\phantom{={}}\color{black}760-(174\times5-296)\\&\color{answerred}=760-(870-296)\\&\color{answerred}=760-574\\&\color{answerred}=186\end{aligned}}
\CalcBlock[方法一]{6}{\begin{aligned}[t]&\phantom{={}}\color{black}900\div[(3+6)\times5]\\&\color{answerred}=900\div[9\times5]\\&\color{answerred}=900\div45\\&\color{answerred}=20\end{aligned}}
\MethodBlock{方法二：除以积，改成连续除}{连续除以两个数，等于除以这两个数的积。}{\begin{aligned}[t]&\phantom{={}}900\div [(3+6)\times 5]\\
&=900\div (3+6)\div 5\\
&=900\div 9\div 5\\
&=100\div 5\\
&=20\end{aligned}}

\par\noindent

\QHeading{第8题　根据运算顺序，在算式中添上小括号或中括号，再计算。}

\Problem{①先算减法，再算除法，最后算乘法\(40\times300-200\div50\)}
\par

\Answer{\Formula{\begin{aligned}[t]&\phantom{={}}40\times[(300-200)\div50]\\&=40\times[100\div50]\\&=40\times2\\&=80\end{aligned}}}
\Problem{②先算除法，再算减法，最后算乘法\(40\times300-200\div50\)}
\Answer{\Formula{\begin{aligned}[t]&\phantom{={}}40\times(300-200\div50)\\&=40\times(300-4)\\&=40\times296\\&=11840\end{aligned}}}
\par\noindent

\QHeading{第9题　解决问题。}

\Problem{①李老师带了300元去买体育用品。买篮球花了120元，剩下的钱买每根10元的跳绳，最多可以买几根？}
\par


\MethodTitle{方法一}
\Answer{\Formula{\begin{aligned}[t]&\phantom{={}}(300-120)\div10\\&=180\div10\\&=18\text{（根）}\end{aligned}}\par 答：最多可以买18根。}

\MethodTitle{方法二：按跳绳单价折算}
\Answer{先看300元能买多少根，再扣去篮球费用相当的根数。}
\Answer{\Formula{\begin{aligned}[t]&\phantom{={}}300\div 10-120\div 10\\
&=30-12\\
&=18\text{（根）}\end{aligned}}}
\Answer{答：最多可以买18根跳绳。}
\Problem{②工程队修一条长800米的路，前4天每天修80米，剩下的计划每天修60米，还要几天才能修完这条路？}
\Answer{\Formula{\begin{aligned}[t]&\phantom{={}}(800-80\times4)\div60\\&=(800-320)\div60\\&=480\div60\\&=8\text{（天）}\end{aligned}}\par 答：还要8天才能修完。}
